\section{Kết luận}
\label{chap:ket-luan}

Trong phép đánh giá nhị phân trên official UNSW-NB15, Stacking-42 và RF-42 đạt Macro-F1 gần nhau nhưng thể hiện cơ cấu lỗi khác nhau: Stacking giảm 142 cảnh báo giả và tăng 118 mẫu Attack bị bỏ sót. CI paired bootstrap của Macro-F1 và MCC chứa 0, vì vậy chưa có bằng chứng về ưu thế tổng thể rõ của Stacking; kết quả cũng không chứng minh hai cấu hình tương đương. Phân tích OOF cho thấy Voting không vượt RF, còn các learner có sai số bổ sung nhưng sự khác biệt đó không tự chuyển thành lợi ích của bộ kết hợp. Kết quả theo họ Attack và novel-vector view giúp mô tả sâu hơn các thay đổi trong benchmark, song chưa xác nhận cơ chế sửa lỗi hay khả năng khái quát sang traffic độc lập.

Đối chiếu RF-17 với RF-42 cho thấy giảm số đặc trưng đi kèm chi phí lịch sử thấp hơn và thay đổi về hiệu năng cùng hai loại lỗi. Vì siêu tham số Random Forest cũng khác, nghiên cứu chưa tách được đóng góp riêng của Feature Selection. Tương tự, dữ liệu hiện có chưa xác định tác động nhân quả của weighting hoặc resampling; và một lần đánh giá cùng bootstrap trên các dự đoán không đo độ ổn định qua các lần huấn luyện. Do đó, các câu trả lời cho RQ1--RQ5 lần lượt xác lập giới hạn của bằng chứng về stability, feature ablation, cơ chế diversity, generalization ngoài split và xử lý mất cân bằng.

Đóng góp thực nghiệm của công trình nằm ở việc đối chiếu các pipeline trên split chính thức bằng nhiều thước đo, trình bày đánh đổi FP/FN cùng độ bất định, phân tích lỗi theo họ tấn công và đặt chi phí quan sát cạnh kết quả dự đoán. Những kết quả không cho thấy ensemble vượt trội là một phần của bằng chứng, đồng thời giới hạn phạm vi kết luận. Tái dựng trên toàn TRAIN trong môi trường phần mềm được xác định hỗ trợ kiểm tra khả năng fit lại cấu hình, nhưng không xác nhận đồng nhất với mô hình lịch sử, không tái xác minh metric TEST và không chứng minh sẵn sàng triển khai.

Nghiên cứu tiếp theo cần đánh giá repeated-training stability, thực hiện ablation Feature Selection theo quy trình nested và matched, kiểm tra riêng weighting/resampling, đồng thời dùng split theo thời gian hoặc nhóm độc lập hay benchmark ngoài. Các bước này sẽ kiểm tra những câu hỏi còn mở bằng dữ liệu và protocol phù hợp, trong khi kết quả của official holdout hiện tại tiếp tục được giữ nguyên.
